%@ pages 421-426
%@ status translated
% Chapter V, section 4: Indecomposable injectives and prime two-sided
% ideals (pp. 421--426). Transcribed from the page images in
% reference/source.pdf, not from the OCR text layer. The section opens
% mid-p. 421 (after Corollary 2 to Theorem 2, which closes section 3)
% and ends on p. 426 with Proposition 8; section 5 begins on the same
% page.
%
% Numbering is automatic (src/preamble.tex). This section prints
% Propositions 6, 7 and 8 (the chapter's run of propositions is at
% Proposition 5 on p. 415), Lemmas 1 and 2 (the lemma counter restarts
% at every section) and Corollaries 1, 2 and 3 (to Proposition 7).
%
% No emendations.

\section{Indecomposable Injectives and Prime Two-Sided Ideals}

--- Let $A$ be a right noetherian ring. If $M$ is an $A$-module,
$\operatorname{Ann} M$ will denote the two-sided ideal formed by the
$a$ such that $M.a = 0$. Likewise, if $m$ is an element of $M$,
$\operatorname{Ann} m$ will denote the right ideal formed by the $a$
such that $m.a = 0$.

Let $I$ be an indecomposable injective $A$-module. If $M$ and $N$ are
two non-zero submodules of $I$, $M \cap N$ is non-zero and one has
\[
\operatorname{Ann}(M \cap N) \supset \operatorname{Ann} M + \operatorname{Ann} N.
\]

The annihilators of the non-zero submodules of $I$ thus form an
increasing filtered family of two-sided ideals of $A$. This family has
a maximal element which we shall denote $A(I)$.

One sees first that $A(I)$ is a prime two-sided ideal; in other words,
$A(I)$ is not equal to $A$ and one has
$a.A.b \subset A(I) \Rightarrow a \in A(I)$ or $b \in A(I)$. This follows
from the following lemma, whose proof is left to the reader:

\begin{lemma}
If $\mathfrak{a}$ is a two-sided ideal of a ring $A$, the following
assertions are equivalent:
\begin{description}
\item[\normalfont a.] $\mathfrak{a}$ is a prime two-sided ideal.
\item[\normalfont b.] There exists a non-zero $A$-module $M$ such that
  $\operatorname{Ann} N = \mathfrak{a}$ for every non-zero submodule
  $N$ of $M$.
\end{description}
\end{lemma}

I also claim that every prime two-sided ideal $\mathfrak{p}$ is of the
form $A(I)$, where $I$ is an indecomposable injective $A$-module:
indeed, let $\bigoplus_k I_k$ be a decomposition of the injective
envelope of $A/\mathfrak{p}$ into indecomposable injectives; let $M_k$
be a submodule of $I_k$ whose annihilator is $A(I_k)$. Then
$M_k \cap A/\mathfrak{p}$ is different from $0$ and has for annihilator
both $\mathfrak{p}$ and $A(I_k)$; whence $\mathfrak{p} = A(I_k)$ for
every $k$, and the result.

It is easy to see that $A/\mathfrak{p}$ is in fact an isotypic
$A$-module (Chapter~IV, \S2), i.e.\ that the injectives $I_k$ are all
isomorphic. Indeed, Corollary~1 to Theorem~2 applies to the ring
$A/\mathfrak{p}$. This ring is therefore contained in a simple ring $B$
which is the injective envelope of $A/\mathfrak{p}$ as a right
$(A/\mathfrak{p})$-module. The injective envelope $I$ of $A/\mathfrak{p}$
considered as an $A$-module contains $B$, and $B$ coincides with the set
of elements of $I$ which are annihilated by $\mathfrak{p}$. The
assertion thus follows from the fact that $B$ is a direct sum of simple
right ideals which are isomorphic co-irreducible $A$-modules.

Let now $M$ be an arbitrary $A$-module, let $I$ be the injective
envelope of $M$ and let $I = \sum_k I_k$ be a decomposition of $I$ as a
direct sum of indecomposable injective modules. The prime two-sided
ideal $\mathfrak{p}$ of $A$ will be said to be \emph{associated} with
$M$ if there exists a $k$ such that $\mathfrak{p} = A(I_k)$. The
indecomposable injectives associated with the same prime two-sided
ideal may then be grouped together in packets
$J_{\mathfrak{p}} = \sum_{A(I_k) = \mathfrak{p}} I_k$. A decomposition
of $I$ is thus obtained, unique up to automorphism, of the form
$I = \sum_{\mathfrak{p}} J_{\mathfrak{p}}$, where $\mathfrak{p}$ runs
through the prime ideals associated with $M$.

In particular, if $M_{\mathfrak{p}}$ is equal to the intersection
$M \cap \left(\sum_{\mathfrak{q} \ne \mathfrak{p}} J_{\mathfrak{q}}\right)$,
one has the equality $0 = \bigcap_{\mathfrak{p}} M_{\mathfrak{p}}$,
where $\mathfrak{p}$ runs through the prime ideals associated with $M$.
Moreover, the decomposition is irredundant and $M/M_{\mathfrak{p}}$ has
$\mathfrak{p}$ as its sole associated prime two-sided ideal. The
following proposition follows from Lemma~1:

\begin{proposition}
The prime two-sided ideal $\mathfrak{p}$ is associated with $M$ if and
only if $M$ possesses a submodule $N$ such that $\mathfrak{p}$ is the
annihilator of $N$ and of every non-zero submodule of $N$.
\end{proposition}

The results of Lesieur-Croisot on the ``tertiary'' decomposition are
thus recovered by a different procedure \cite{lesieur-croisot-2}.

We now propose to see under which conditions the theory obtained is no
finer than that of Lesieur-Croisot, i.e.\ when the correspondence
established between indecomposable injectives and prime two-sided ideals
is bijective. For this, we shall assume hypothesis (H) satisfied:

(H) If $I$ is a right ideal of $A$ and if $\mathfrak{p}$ is the
annihilator of $A/I$, there exists a finite number of elements
$x_1, \ldots, x_r$ of $A/I$ such that
\[
\mathfrak{p} = \bigcap_{1 \le i \le r} \operatorname{Ann} x_i.
\]

If condition (H) is satisfied, I claim that, for every indecomposable
injective $I$, the injective envelope of $A/A(I)$ is a direct sum of a
finite number of modules isomorphic to $I$: let indeed $M$ be a
submodule of $I$ whose annihilator is $A(I)$. One manifestly has
\[
A(I) = \bigcap_{\substack{x \ne 0 \\ x \in M}} \operatorname{Ann} x.
\]

But it follows from (H) that $A(I)$ is the intersection of a finite
number of $\operatorname{Ann} x$, say
$\operatorname{Ann} x_1, \ldots, \operatorname{Ann} x_r$; there is
therefore a monomorphism from $A/A(I)$ into the finite direct sum
$\bigoplus_{i=1}^{r} x_i.A$. The injective envelope of this sum is
isotypic; whence the result:

\begin{lemma}
If condition (H) is satisfied, the correspondence $I \leadsto A(I)$
between types of indecomposable injectives and prime two-sided ideals
is bijective.
\end{lemma}

It remains to give examples where condition (H) is satisfied:

\begin{description}
\item[\normalfont a.] Every right ideal of $A$ is two-sided.
\item[\normalfont b.] The ring $A$ is right artinian. The radical
  $\mathfrak{r}(A)$ of $A$ is then nilpotent and every prime two-sided
  ideal contains $\mathfrak{r}(A)$. The prime two-sided ideals of $A$
  are therefore in bijective correspondence with those of the
  semisimple ring $A/\mathfrak{r}(A)$, i.e.\ with the irreducible
  representations of $A$.
\item[\normalfont c.] The centre $Z(A)$ of $A$ is a noetherian ring and
  $A$ is a $Z(A)$-module of finite type: let $M$ be an $A$-module of
  finite type, $x$ and $y$ two elements of $M$. If $x'$ and $y'$ denote
  the sets formed by the elements of $M$ which are annihilated by
  $\operatorname{Ann} x$ and $\operatorname{Ann} y$, one sees that $x'$
  and $y'$ are modules of finite type over $Z(A)$. Moreover, one has
  the equality
  \[
  \operatorname{Ann}(x' + y') = \operatorname{Ann} x \cap \operatorname{Ann} y.
  \]
  Since $M$ is a noetherian $Z(A)$-module, there exists a finite number
  of elements $x_1, \ldots, x_r$ of $M$ such that
  \[
  x_1' + \cdots + x_r' = M.
  \]
  Condition (H) is therefore satisfied.
\item[\normalfont d.] A counterexample: let $k$ be a commutative field
  of characteristic $0$ and $k(X)$ the field of rational fractions in
  one indeterminate $X$ over $k$. Let $d$ be the $k$-derivation
  $P(X) \mapsto P'(X)$ of $k(X)$ and $A$ the ring of operators generated
  by $d$ and the homotheties of $k(X)$. The ring $A$ is a domain, every
  right ideal is principal and there is no two-sided ideal distinct from
  $0$ or $A$ (\cite{bourbaki-algebre}, \S5, exerc.~13). In particular,
  the prime two-sided ideals associated with a simple $A$-module and
  with $A_d$ are zero. Since these modules do not have the same
  injective envelope, one sees that the correspondence
  $I \leadsto A(I)$ is not bijective.
\end{description}

The set of prime two-sided ideals of $A$ will henceforth be called the
\emph{prime spectrum} of $A$ [notation: $\operatorname{Spec}(A)$]. We
know that it is possible to associate with every prime two-sided ideal
$\mathfrak{p}$ an indecomposable injective $A$-module $I_{\mathfrak{p}}$
satisfying the following condition: the injective envelope of
$A/\mathfrak{p}$ is a finite direct sum of modules isomorphic to
$I_{\mathfrak{p}}$. The map $\mathfrak{p} \leadsto I_{\mathfrak{p}}$
defines an injection of $\operatorname{Spec}(A)$ into the spectrum of
the category $\operatorname{mod} A$. We shall always identify
$\operatorname{Spec}(A)$ with the image of this injection.

\begin{proposition}
Let $A$ be a right noetherian ring and $\mathbf{C}$ a localizing
subcategory of $\operatorname{mod} A$. A prime two-sided ideal
$\mathfrak{p}$ belongs to $\operatorname{Sp}(\mathbf{C})$
(Chapter~IV, \S1) if and only if $A/\mathfrak{p}$ is an object of
$\mathbf{C}$.
\end{proposition}

It is clear that $\mathfrak{p}$ belongs to
$\operatorname{Sp}(\mathbf{C})$ if $A/\mathfrak{p}$ is an object of
$\mathbf{C}$. Suppose conversely that $\mathfrak{p}$ belongs to
$\operatorname{Sp}(\mathbf{C})$. This means that $I_{\mathfrak{p}}$
contains a non-zero submodule belonging to $\mathbf{C}$. It follows that
$A/\mathfrak{p}$ contains a non-zero submodule $N$ belonging to
$\mathbf{C}$. We denote by $S$ the multiplicative set formed by the
regular elements of $A/\mathfrak{p}$ (cf.~\S3). If $P$ is a subgroup of
$N$, the right ideal $\operatorname{Ann} P$ of $A/\mathfrak{p}$ formed by
the $a$ such that $P.a$ is zero is the intersection of $A/\mathfrak{p}$
with a right ideal of the simple ring $(A/\mathfrak{p})_S$. It follows
that the right ideals of $A$ of the form $\operatorname{Ann} P$ satisfy
the descending chain condition. In particular, there exists a finite
number of elements $x_1, \ldots, x_r$ of $N$ such that
\[
0 = \bigcap_{i=1}^{r} \operatorname{Ann} x_i.
\]

There is therefore an $A$-linear map from $A/\mathfrak{p}$ to the direct
sum of the modules $x_i.A$. It follows that $A/\mathfrak{p}$ belongs to
$\mathbf{C}$. \hfill Q.E.D.

If $s$ is an element of $\operatorname{Sp}(\operatorname{mod} A)$, and if
$I_s$ is an indecomposable injective $A$-module of type $s$, the
two-sided ideal $A(I_s)$ belongs to $\operatorname{Spec}(A)$. We shall
say that the map $s \leadsto A(I_s)$ is the \emph{canonical projection}
from $\operatorname{Sp}(\operatorname{mod} A)$ onto
$\operatorname{Spec}(A)$. This projection is the identity on
$\operatorname{Spec}(A)$; it is therefore surjective. When it is
bijective, we say that the ring $A$ \emph{has enough two-sided ideals}.

Let $A$ be a ring, right noetherian and having enough two-sided ideals.
If $F$ is a topologising and idempotent set of right ideals, we have
seen that the product $I.\mathfrak{m}$ of two right ideals $I$ and
$\mathfrak{m}$ belongs to $F$ provided $I$ and $\mathfrak{m}$ belong to
$F$. The multiplicative law of $A$ therefore equips the set $F$ with a
monoid structure. I claim that the submonoid of $F$ generated by the
prime two-sided ideals belonging to $F$ is cofinal: the subcategory
$\mathbf{F}$ is indeed generated by the modules $A/\mathfrak{p}$, where
$\mathfrak{p}$ runs through the prime two-sided ideals belonging to $F$
(Proposition~7 and Corollary~2 to Proposition~2, Chapter~IV). Every
noetherian module $M$ belonging to $\mathbf{F}$ consequently possesses a
composition series
\[
0 = M_0 \subset M_1 \subset \cdots \subset M_r \subset M_{r+1} = M,
\]
whose quotients $M_{i+1}/M_i$ are annihilated by a prime two-sided ideal
$\mathfrak{p}_i \in F$. One deduces from this that the product
$\mathfrak{p}_r.\mathfrak{p}_{r-1}\ldots\mathfrak{p}_0$ annihilates $M$.
If $M$ is the quotient of $A$ by an element $I$ of $F$, it follows that
$I$ contains the product
$\mathfrak{p}_r.\mathfrak{p}_{r-1}\ldots\mathfrak{p}_0$. Proposition~7
therefore implies the

\begin{corollary}
Let $A$ be a right noetherian ring having enough two-sided ideals and
$E$ a set of prime two-sided ideals of $A$. The set $F$ of right ideals
containing a product of the form
$\mathfrak{p}_r.\mathfrak{p}_{r-1}\ldots\mathfrak{p}_0$, where
$\mathfrak{p}_i \in E$, is topologising and idempotent. Conversely,
every topologising and idempotent set of right ideals is of this type.
\end{corollary}

If the ring $A$ has enough two-sided ideals, the prime spectrum of $A$
coincides with the spectrum of the category $\operatorname{mod} A$. The
latter is thus equipped with a noetherian ordered set structure.
Moreover, for every localizing subcategory $\mathbf{C}$ of
$\operatorname{mod} A$, $\operatorname{Sp}(\mathbf{C})$ coincides with the
set of prime two-sided ideals $\mathfrak{p}$ such that $A/\mathfrak{p}$
belongs to $\mathbf{C}$. If $\alpha$ is an ordinal, we denote in
particular by $E_\alpha$ the set of prime ideals such that
$\operatorname{Kdim} A/\mathfrak{p} \le \alpha$ (see Chapter~IV, \S1).
One then has the

\begin{corollary}
Let $A$ be a right noetherian ring having enough two-sided ideals. With
the notation above, the sets $E_\alpha$ may be defined as follows:
\begin{itemize}
\item $E_{-1}$ is empty.
\item If the ordinal $\alpha$ has a predecessor $\beta$, $E_\alpha$ is
  formed by the prime two-sided ideals $\mathfrak{p}$ such that every
  prime two-sided ideal containing $\mathfrak{p}$ and distinct from
  $\mathfrak{p}$ belongs to $E_\beta$.
\item If $\alpha$ is a limit ordinal, $E_\alpha$ is the union of the
  $E_\beta$ for $\beta < \alpha$.
\end{itemize}
\end{corollary}

It is clear that $E_{-1}$ is empty. Let therefore $\gamma$ be an
arbitrary ordinal and suppose that the statements of Corollary~2 hold
for every ordinal $\alpha < \gamma$. If $\gamma$ is a limit ordinal,
these statements manifestly remain true when $\alpha$ is equal to
$\gamma$. Let us examine the case where $\gamma$ has a predecessor
$\beta$, and let $T_\beta$ be the canonical functor from
$\operatorname{mod} A = \mathbf{A}$ to
$\mathbf{A}/\mathbf{A}_\beta$ (the notation being that of
Chapter~IV):

Let $\mathfrak{p}$ be a prime two-sided ideal such that
$\operatorname{Kdim} A/\mathfrak{p} = \gamma$. On the other hand, let
$\mathfrak{q}$ be a prime two-sided ideal containing $\mathfrak{p}$ and
distinct from $\mathfrak{p}$. It is clear that
$\operatorname{Kdim} A/\mathfrak{q}$ is less than or equal to
$\operatorname{Kdim} A/\mathfrak{p}$; if equality held, the
Jordan-H\"older quotients of $T_\beta(A/\mathfrak{p})$ would not all be
isomorphic. There would follow the existence of a non-zero submodule $N$
of $A/\mathfrak{p}$ which would not generate the same localizing
subcategory of $\operatorname{mod} A$ as $A/\mathfrak{p}$. Since there
exists an $A$-linear injection of $A/\mathfrak{p}$ into a direct sum of
modules isomorphic to $N$ (cf.~the proof of Proposition~7), this is
impossible. One has therefore the inequality
$\operatorname{Kdim} A/\mathfrak{q} < \operatorname{Kdim} A/\mathfrak{p}$.

Conversely, suppose that $A/\mathfrak{q}$ belongs to
$\mathbf{A}_\beta$ for every prime ideal $\mathfrak{q}$ containing
$\mathfrak{p}$ and distinct from $\mathfrak{p}$. Let $I$ be a right
ideal of $A/\mathfrak{p}$ which is maximal among the right ideals
$\mathfrak{m}$ such that $A/\mathfrak{m}$ does not belong to
$\mathbf{A}_\beta$. Then $T_\beta(A/I)$ is a simple object; it follows
that there exists a two-sided prime ideal associated with $A/I$. This
prime ideal contains $\mathfrak{p}$, and the hypothesis we have made
shows that it coincides with $\mathfrak{p}$. Since $A/I$ belongs to
$\mathbf{A}_\gamma$, $\operatorname{Sp}(\mathbf{A}_\gamma)$ contains
$\mathfrak{p}$. It follows that $A/\mathfrak{p}$ belongs to
$\mathbf{A}_\gamma$ (Proposition~7).

\begin{corollary}
Let $A$ be a right noetherian ring having enough two-sided ideals,
$\mathfrak{p}$ a prime two-sided ideal of $A$ and $\eta$ a finite
ordinal. The following assertions are equivalent:
\begin{description}
\item[\normalfont a.] $\operatorname{Kdim} A/\mathfrak{p} \le \eta$.
\item[\normalfont b.] Every chain of prime two-sided ideals containing
  $\mathfrak{p}$ has at most $\eta + 1$ elements.
\end{description}
\end{corollary}

Corollary~3 follows directly from Corollary~2. If $A$ is a commutative
noetherian ring, it shows that the notion of Krull dimension we
introduced in $\operatorname{mod} A$ coincides with the classical notion.

To conclude this section, we recall the examples we have given:

\begin{proposition}
A right noetherian ring $A$ has enough two-sided ideals if one of the
following conditions holds:
\begin{description}
\item[\normalfont a.] $A$ is right artinian.
\item[\normalfont b.] Every right ideal is two-sided.
\item[\normalfont c.] The centre $Z(A)$ of $A$ is a noetherian ring and
  $A$ is a $Z(A)$-module of finite type.
\end{description}
\end{proposition}
